\begin{minipage}{0.65\linewidth}
    On applique le théorème d'énergie cinétique~:
    
    \[
        \Delta E_c=\sum W_{AB}(\vec{F}_{\text{ext}})
    \]
    avec~:
    \[
        \left\{
        \begin{array}{l}
            \Delta E_c=\frac{1}{2}m v_B^2-\frac{1}{2}m v_A^2=
            \frac{1}{2}m v^2-0 \\[3mm]
            \sum W_{AB}(\vec{F}_{\text{ext}})=W(\vec{P})=mg(z_{A}-z_{B})=mgh
        \end{array}\right.
    \]
    donc~:
    \[
        \frac{1}{2}m v^2=mgh
        \quad\Rightarrow\
        \mathbox{v=\sqrt{2gh}}
    \]

    Application numérique~:
    \[
        v=\sqrt{2\times 9,8\times 300}
        \quad\Rightarrow\
        v\approx \qty{77}{\v}
    \]
    
\end{minipage}
\hfill
\begin{minipage}{0.34\linewidth}

\begin{figure}[H]
    \centering
    \begin{tikzpicture}
        \def\h{4}
        \Sol{3cm}{3mm}{sol}
        \coordinate (O) at (sol.north west);
        \draw[-{Latex},thick] (O) +(0,-0.5) -- +(0,5) node[right] {$z$};
        \foreach \z/\n in {\h+0.15/A,0.15/B}{
            \draw[thick] ($(O)+(-2pt,\z)$) node[left] {$z_{\n}$} -- ++(4pt,0);
            \node[circle,ball color=black!40,inner sep=0pt,
                minimum size=3mm,label=$\n$] (\n) at ($(O)+(2,\z)$) {};
            \node[circle,inner sep=0.6pt,fill] (\n) at (\n.center) {};
            \draw[dashed] (\n) -- (\n -| O);
        }
        \draw[->,thick,blue] (A.center) -- ++(0,-1.2) node[left] {$\vec{P}$};
        \draw[<->,shorten <=1pt,shorten >=1pt] ([xshift=0.5cm,yshift=0.15cm]O) --
            node[right] {$h$} ++(0,\h);
        \node[right=2mm] at (A.east) {$v_{A}=0$};
        \node[right=2mm] at (B.east) {$v_{B}=v$};
    \end{tikzpicture}
\end{figure}
\end{minipage}
